\subsection{Paired comparisons and the radiance--reflectance trade-off}\label{sec:paired-tails}

Paired differences compare two families on the same split. Table~\ref{tab:paired} reports them
for the main and wide pipelines, with positive gains favoring the candidate. The wide feature
kernel improves on the input-scaled kernel at all ten splits, by $0.0165\pm0.0066$ radiance
percentage points. The wide stack's improvement over its own feature kernel is $0.0006\pm0.0026$
points, with eight wins and two losses: small against the split-to-split variation of the
difference itself.

\begin{table}[t]
\centering\small
\resizebox{\textwidth}{!}{\begin{tabular}{lrrr}
\toprule
candidate / reference & mean gain $\pm$ SD [pp] & gain range [pp] & wins/10 \\
\midrule
Input-scaled / isotropic KRR & 0.1436$\pm$0.0122 & [0.1196, 0.1585] & 10 \\
Residual correction / network & 0.2464$\pm$0.0090 & [0.2365, 0.2615] & 10 \\
Narrow stack / input-scaled KRR & 0.0083$\pm$0.0066 & [-0.0054, 0.0185] & 9 \\
Wide feature kernel / input-scaled KRR & 0.0165$\pm$0.0066 & [0.0054, 0.0295] & 10 \\
Wide / narrow feature kernel & 0.0191$\pm$0.0079 & [0.0096, 0.0352] & 10 \\
Wide stack / wide feature kernel & 0.0006$\pm$0.0026 & [-0.0053, 0.0037] & 8 \\
\bottomrule
\end{tabular}
}
\caption{Paired radiance contrasts on seeds 101--110. Gain is the reference error minus the
candidate error, in percentage points (pp). SD is the sample standard deviation of the ten
paired differences; the range and the win count describe the same ten differences.}
\label{tab:paired}
\end{table}

Radiance and reflectance order the families differently. Table~\ref{tab:tails} reports the
all-band reflectance tails beside the medians; each entry averages a within-split statistic, so
the mean of ten 95th percentiles is not a pooled 95th percentile. For the width-512 pipeline the
convex stack lowers the radiance error relative to the feature kernel but raises the mean 95th
percentile of the reflectance error from $7.75$ to $29.44$ points, and the tail is worse at all
ten splits, by factors between $2.63$ and $4.94$. The width-2000 pipeline repeats the reversal:
$7.57$ points for the feature kernel and $22.71$ for the stack, worse at all ten splits by
factors between $2.36$ and $3.79$. The near-tie at $0.079\%$ against $0.078\%$ radiance error
hides a threefold difference in the inverse.

\begin{table}[t]
\centering\small
\resizebox{\textwidth}{!}{\begin{tabular}{lrrrr}
\toprule
model & radiance [\%] & median [pp] & 95th percentile [pp] & RMSE [pp] \\
\midrule
Cubic ridge & 0.760$\pm$0.011 & 0.471$\pm$0.010 & 70.68$\pm$0.02 & 4.18e+03 \\
Isotropic KRR & 0.239$\pm$0.016 & 0.141$\pm$0.018 & 69.53$\pm$0.42 & 4.18e+03 \\
Input-scaled KRR & 0.095$\pm$0.009 & 0.036$\pm$0.004 & 31.16$\pm$6.21 & 1.34e+04 \\
Narrow feature kernel & 0.098$\pm$0.007 & 0.054$\pm$0.004 & 7.75$\pm$1.13 & 1.78e+03 \\
Narrow convex stack & 0.087$\pm$0.006 & 0.048$\pm$0.003 & 29.44$\pm$5.89 & 1.25e+03 \\
Wide feature kernel & 0.079$\pm$0.009 & 0.045$\pm$0.006 & 7.57$\pm$0.95 & 1.40e+03 \\
Wide convex stack & 0.078$\pm$0.011 & 0.043$\pm$0.005 & 22.71$\pm$4.42 & 5.81e+03 \\
\bottomrule
\end{tabular}
}
\caption{All-band retrieval metrics over the ten splits of the main, concatenated-feature and
wide pipelines. Radiance is the mean relative $L^2$ error; the reflectance median, 95th
percentile and RMSE are in percentage points at $\rho=0.7$. The first three columns show mean
$\pm$ sample standard deviation across splits and RMSE its mean. The inverse is neither
clipped nor restricted to the physical domain; every inverse in these runs is finite, and the
RMSE is dominated by large finite values near zeros of the denominator.
Section~\ref{sec:domain} gives the same comparison on the physical domain.}
\label{tab:tails}
\end{table}

The input-scaled kernel shows the same split between center and tail: its median error is only
$0.0365$ points but its mean 95th percentile is $31.16$ points, while the feature kernels have
slightly larger medians and much smaller tails. The kernel on concatenated features has the
smallest tail of all, $5.93$ points. The identities of Section~\ref{sec:stability} tie the
inverse error to the predicted denominator as well as to the forward error, which is why
optimizing a component or radiance error alone need not control the reflectance tail.
